\section{Specialization to the finite field}
\label{sec:final}

We now have a perverse eigensheaf with an action of $\varphi^m$ on
the lifted pointed curve over $\Dbar$.  The arithmetic torsor tower
of Proposition~\ref{prop:lifted-parameter} identifies its eigenvalue
after mixed-characteristic specialization with the original
arithmetic parameter, including its Frobenius structure.

\begin{proof}[Proof of Theorem~\ref{thm:main}]
Take $m$ and $M_{\Dbar}$ from
Proposition~\ref{prop:charzero-equivariant}.  Recall that
$R=W(k)$, $K=\overline{\Frac(R)}$, and $\Dbar\subset K$ is
preserved by $\varphi$.  Pullback under this algebraically closed
field extension gives
\[
 M_K\quad\text{on}\quad
 \Bun_{G,B,\mathfrak x}(\mathscr X_K).
\]
It remains nonzero, locally constructible, and perverse, and
inherits the $\varphi^m$-structure.  These assertions are checked
on finite-type smooth charts: extension of algebraically closed
fields preserves nonzero constructible stalks, support dimensions,
and Verdier duality, hence perversity.  Base change for the bounded
Hecke correspondences, with the tensor Satake comparison of
Section~\ref{sec:conventions}, transports the entire eigenstructure.
Its eigenvalue is the base change of $\sigma_1$, with the action
provided by the arithmetic tower.

Consider the relative ordinary-level bundle stack
\begin{equation}\label{eq:final-relative-stack}
 \mathscr A=
 \Bun_{G,B,\mathfrak x}(\mathscr X_R/R),
 \qquad L=\mathscr U_R,
 \qquad M=\Psi M_K\quad\text{on }\mathscr A_k.
\end{equation}
We verify the hypotheses of
Proposition~\ref{prop:equivariant-specialization} and
Lemma~\ref{lem:specialization-detection} in this second family.
The ring $R$ is an excellent complete strictly henselian DVR with
algebraically closed residue field $k$, and $\ell$ is invertible
because $\ell\ne p$.  The stack $\mathscr A$ is smooth and locally
of finite type over $R$: $H^2$ of the adjoint bundle vanishes on
the fibers of the smooth relative curve, and adding a flag is a
smooth $G/B$-fibration.  The leg space $L$ is smooth of relative
dimension one, with special fiber $U$.

Since all legs avoid $\mathfrak x$, their spherical Hecke
correspondences have the ordinary split affine Grassmannian
models in frames.  Bounded output maps are representable and
proper, both exact-position maps over bundle-and-leg are smooth,
and the open-stratum kernels have the fixed relative Satake
normalization.  The special-fiber cotangent and
transverse-modification geometry is the ordinary Borel-level
case of \cite[Section~3]{CT}.  Its group hypotheses H1--H4 hold
by Lemma~\ref{lem:lie-hypotheses}, using $p>n$.
Finally, the invariant lattices supplied by the arithmetic tower
have lisse finite reductions on $\mathscr U_R$, with their
$\varphi$-actions.  Their specializations, including all tensor
maps, are exactly the local systems associated to $\rho$.
Thus all the geometric and eigenvalue hypotheses of both cited
specialization statements hold.

It follows that $M$ is locally constructible and perverse, and
has the natural multi-leg eigenisomorphisms
\begin{equation}\label{eq:final-eigen}
 \Hecke_{I,(V_i)}(M)
 \simeq M\boxtimes
                 \mathop{\boxtimes}_{i\in I}(V_i)_{\rho|_{\pi_1(U)}}.
\end{equation}
The same proposition gives a $\varphi^m$-structure on $M$ and
equivariance of these maps.  The special-fiber action of
$\varphi^m$ is $q^m$-Frobenius.  On the eigenvalue system its
action is precisely the one obtained by restricting
$\rho$ to $\pi_1(U_{0,\F_{q^m}})$, because the finite torsors used
for the lift descended from that arithmetic representation.
This identifies the arithmetic eigenvalue, as well as its
geometric restriction.

We still have to prove $M\ne0$.  Choose a $K$-point $(P,\beta)$
with nonzero stalk of $M_K$, where $\beta$ is its Borel reduction
at $\mathfrak x_K$.  Use the trivial $G$-bundle on $\mathscr X_R$
as reference, and lift a point of $U(k)$ to a section
$y\in\mathscr U_R(R)$ by smoothness and the henselian property.
The affine-curve argument of Section~\ref{sec:nodal} trivializes
$P$ off $y_K$, compatibly with its determinant.  The resulting
modification of the reference bundle lies in a finite Schubert
bound, whose Hecke space is proper over the reference input and
$y$.  After a finite trait extension it extends to a bundle
$\mathscr P$ on the whole relative curve.

The remaining datum is the flag.  Its generic value $\beta$ is a
point of the associated flag bundle
$\mathscr P|_{\mathfrak x}\times^G G/B$ over the extended trait.
This bundle is proper, so the valuative criterion extends
$\beta$ as well.  We have thereby extended the point
$(P,\beta)$ to a section of $\mathscr A$.  Lift this section to a
smooth finite-type chart through its special value, using
henselian lifting as in the nodal argument.  The section has a
nonzero generic stalk, so the closure of the nonzero-stalk locus
on that chart meets its special fiber.  All hypotheses having
been verified above, Lemma~\ref{lem:specialization-detection}
proves $M\ne0$.

Apply now the field-level nilpotent-detection theorem
\cite[Theorem~4.1]{CT}.  For the level stacks over
$k=\overline{\F}_q$ under its Lie-theoretic hypotheses, it states
that a locally bounded constructible
complex with spherical eigenisomorphisms and lisse eigenvalues
has singular support in the cone of generically nilpotent Higgs
fields on every smooth finite-type chart.  Here $\mathscr A_k$
is the ordinary Borel-level stack, the characteristic assumptions
hold, and \eqref{eq:final-eigen} supplies those lisse eigenvalues.
At an ordinary flag $\beta$, the cotangent residue condition is
\[
 \operatorname{Res}_{x}\theta
        \in\operatorname{Lie}R_u(B_\beta).
\]
The cone in the detection theorem is therefore exactly
$\Lambda_{\mathrm{par}}$ from Theorem~\ref{thm:main}.  This proves
$\SS(M)\subset\Lambda_{\mathrm{par}}$ in the stated chartwise sense.

Equip $M$ with its $\varphi^m$-structure and denote the resulting
Weil sheaf over $\F_{q^m}$ by $M_0$.  The split Weil Satake
comparison of Section~\ref{sec:conventions} identifies the
transported Hecke functors with the relative-Satake-normalized
Weil functors in the theorem.  Thus \eqref{eq:final-eigen} is an
isomorphism of Weil sheaves for every finite leg set and every
collection of representations.

Every operation above transports the full eigenstructure by
comparisons natural for the semilinear generator.  These
comparisons were constructed on the successive-modification
correspondences and on every collision diagonal, with the tensor
unit and permutations respected throughout.  The resulting
unit, convolution, permutation, and fusion diagrams therefore
commute as Weil diagrams, including all iterated collisions.
A structure for the group generated by Frobenius is all that is
required for a Weil sheaf; no profinite descent assertion is
used.  The resulting $M_0$ has all the claimed properties.
\end{proof}
